% !TEX root =  Lecture7.tex


%%%%%%%%%%%%%%%%%%%%%%%%%%%%%%%%%%%%
% Recap/Agenda
%%%%%%%%%%%%%%%%%%%%%%%%%%%%%%%%%%%%
\begin{frame}
\frametitle{Module 1: Modeling and Linear Regression}
    \begin{itemize}
    \item \hl{Previously:} logistic regression --- modeling a binary response, odds ratios, and maximum likelihood.
    \item \hl{This time:} choosing between models.
      \begin{itemize}
      \item multiple logistic regression and its candidate models
      \item model selection with adjusted $R^2$ and \hl{AIC}
      \item multicollinearity: what goes wrong when predictors overlap
      \end{itemize}
    \item \hl{Reading:} IMS Chapters 8.4 and 9.3.
    \item \hl{Homework for today:}
    \item \hl{Homework for next class:}
    \end{itemize}
\end{frame}


%%%%%%%%%%%%%%%%%%%%%%%%%%%%%%%%%%%%
% Learning objectives:
%%%%%%%%%%%%%%%%%%%%%%%%%%%%%%%%%%%%
\begin{frame}
    \frametitle{Lecture Objectives}
    \goals{%
    \begin{itemize}
    \item explain why $R^2$ can never choose a model, and what adjusted $R^2$ fixes;
    \item compute and interpret \hl{AIC}, and use it to compare candidate models;
    \item describe \hl{multicollinearity} and how it distorts coefficients and standard errors;
    \item select a model that balances goodness of fit against complexity.
    \end{itemize}}
    \objectives{%
    \begin{itemize}
    \item \textbf{(M1, LO4)} Model Selection and $R^{2}$ and Adjusted $R^{2}$
    \item \textbf{(M1, LO3)} Using Regression Models --- multicollinearity
    \end{itemize}}
\end{frame}
